\section{Results}

\subsection{Typed priors for unknown feasibility}
In N4-A, type-conditioned priors are the only difference between `ORION_TYPED_VOI` and the uniform-prior VOI ablation. Over 300 paired episodes, the typed arm has mean utility 3.291 versus 2.180 for uniform VOI and 4.612 for the full oracle, recovering about 71\% of oracle utility. Blind optimistic commitment is driven to $-13.619$ mean utility and succeeds in only 19\% of episodes, satisfying the hostile-world validity gate.

\subsection{Scope-bound failure-receipt reopening}
N4-B separates changes that invalidate a failure receipt from irrelevant context churn. Pooled mean utility is 3.199 for scoped reopening, 2.782 for never reopening, $-7.813$ for unscoped-change reopening, and $-9.225$ for always reopening. In the deliberately wasteful regime, always reopening collapses to about $-13.406$ while never reopening remains positive; the scoped mechanism retains the no-giveback property required by the protocol.

The scoped-versus-never gap is not a supported result, and we state this explicitly because the pooled means invite the opposite reading. Under the paired analysis of Table~\ref{tab:N4-U1}, scoped reopening is not separated from never reopening in either regime: the paired mean difference is $0.060$ with 95\,\% interval $[-0.540, 0.634]$ in the wasteful regime and $0.774$ with interval $[-0.663, 2.254]$ in the stale regime, on 200 pairs each. Both intervals contain zero. The wasteful-regime comparison is additionally 79\,\% ties, so the pooled means are separated mostly by a small minority of discordant episodes. What N4-B does establish at this sample size is the scoped-versus-\emph{unscoped} contrast, where both intervals exclude zero by a wide margin ($[13.624, 16.451]$ and $[5.740, 8.255]$). The mechanism claim this study supports is therefore that scoping reopening beats reopening on unscoped change --- not that reopening under scope beats never reopening at all.

\subsection{Targeted verification under interval costs}
N4-C gives every active arm the same budget of four edge verifications. Interval-dominance targeting yields mean scalarized regret 0.1096, compared with 0.2518 for random verification, 0.2621 for midpoint optimization without targeted verification, 0.7755 for worst-case optimization, and 1.2679 for best-case optimization. The benefit is therefore attributed to where verification is spent, not verification volume.

\subsection{Full-chain certificate transport}
N4-D contains 200 honest and 200 laundering chains. The full-chain checker has recall 1.000 and false-positive rate 0.000. It detects all missing-receipt, spoofed-summary, and 68 deep-splice attacks. Last-hop inspection has zero recall on the deep-splice class. This is a construction-level provenance result; the identifiers are seeded values rather than cryptographic digests.

\subsection{Decision-coupled active experiments}
In N4-E, pure information gain is attracted by high-entropy facts that are irrelevant to the downstream choice. It spends 36.6\% of its probes on decoys and obtains mean utility 7.121, while the decision-coupled selector spends none on decoys and reaches 9.266 under the shared stopping rule. The decoy-attraction gate is part of world validity.

\subsection{Remint and transport across edits}
N4-F3 compares typed remint/transport with naive carry-forward and matched-budget re-derivation. In the stale-hostile regime, naive carry-forward fails on 99.5\% of episodes and has mean utility about $-15.916$; typed transport reaches 9.421 versus 7.157 for re-derivation. In the prespecified remint-unnecessary regime, all relevant arms tie exactly at 11.809659685355605 and the typed arm spends zero remints. A positive advantage there would have invalidated the study.

\subsection{Paired uncertainty for every comparison}
The means reported above are point estimates. Table~\ref{tab:N4-U1} gives the paired 95\,\% percentile bootstrap interval, the pair count, and the win/tie/loss decomposition for each of the twelve N4 comparisons, transcribed from the committed publication analysis. Three readings matter for how far these results should be taken.

First, two comparisons are undetermined: both N4-B scoped-versus-never contrasts have intervals containing zero, as discussed above. Every other non-degenerate comparison excludes zero.

Second, several intervals that exclude zero rest on a minority of discordant pairs, and the tie column is the honest measure of how often the mechanism changes anything at all. N4-A typed-versus-uniform VOI wins 29\,\% of pairs against 13\,\% losses with 58\,\% ties; N4-C wins 26\,\% against 7\,\% with 67\,\% ties. Both differences are consistent in direction, but they are carried by roughly a third of episodes, not by a uniform shift. The large-margin results --- N4-E versus information gain (95\,\% wins) and N4-F3 versus naive carry-forward (69\,\% wins, zero losses) --- are the ones that hold on most episodes.

Third, the intervals are per-comparison and unadjusted for multiplicity. Twelve comparisons are reported, nine of which exclude zero. Any family-wise adjustment widens all of them. The one direction that follows without further computation is that the two undetermined contrasts can only remain undetermined. We do not assert here which of the nine would survive adjustment. The committed artifact stores summary statistics rather than the paired difference vectors, so corrected \emph{bootstrap intervals on the mean} cannot be derived from it. A family-wise-correctable exact test is available by a different route, over the win/loss counts rather than the means, and is reported separately; we keep the unadjusted intervals here because the mean is the registered estimand for these comparisons. Comparisons whose interval lies close to zero relative to its width --- N4-C and N4-E against the LLM proxy --- are the ones a family-wise correction would put at risk.

% Generated by papers/orion-08-typed-state/scripts/emit_paired_uncertainty_table.py
% Source: PUBLICATION_PAIRED_ANALYSIS_V1.json (schema ORION.Q4.PublicationPairedAnalysis.v1). Do not edit by hand.
% Every value is transcribed from the source JSON. No new estimate is derived.
\begin{table}[t]
  \centering
  \caption{Paired uncertainty for every N4 comparison in the committed publication analysis. Intervals are 95\,\% percentile bootstrap intervals over paired episode differences (5000 deterministic resamples, seeds recorded per row in the source artifact). Win/tie/loss is the paired decomposition over the same $n$ pairs. Intervals are per-comparison and are \emph{not} adjusted for the twelve comparisons shown; any family-wise adjustment widens them, so rows already containing zero would continue to contain zero. Rows marked $\dagger$ have an interval containing zero: the contrast is \emph{undetermined} at this $n$ and is not a positive result. The row marked $\ddagger$ is the prespecified remint-unnecessary regime, where the exact tie on all pairs is the passing outcome and a positive advantage would have invalidated the study; its interval is $[0,0]$ by construction and is not an undetermined contrast.}
  \label{tab:N4-U1}
  \footnotesize
  \begin{tabularx}{\textwidth}{l>{\raggedright\arraybackslash}X r r r r}
    \toprule
    Study & Comparison & Mean diff. & 95\,\% CI & $n$ & win/tie/loss \\
    \midrule
    N4-A & typed VOI vs.\ greedy known-graph & 2.933 & $[2.556,\ 3.301]$ & 300 & 0.55/0.30/0.15 \\
    N4-A & typed VOI vs.\ uniform-prior VOI & 1.111 & $[0.833,\ 1.400]$ & 300 & 0.29/0.58/0.13 \\
    N4-B & scoped vs.\ never reopen (wasteful)$\dagger$ & 0.060 & $[-0.540,\ 0.634]$ & 200 & 0.17/0.79/0.04 \\
    N4-B & scoped vs.\ unscoped reopen (wasteful) & 15.050 & $[13.624,\ 16.451]$ & 200 & 0.78/0.11/0.11 \\
    N4-B & scoped vs.\ never reopen (stale)$\dagger$ & 0.774 & $[-0.663,\ 2.254]$ & 200 & 0.59/0.20/0.21 \\
    N4-B & scoped vs.\ unscoped reopen (stale) & 6.973 & $[5.740,\ 8.255]$ & 200 & 0.52/0.35/0.13 \\
    N4-C & interval-Pareto vs.\ random verification & 0.142 & $[0.100,\ 0.187]$ & 400 & 0.26/0.67/0.07 \\
    N4-E & decision VOI vs.\ information gain & 2.146 & $[1.976,\ 2.299]$ & 400 & 0.95/0.03/0.02 \\
    N4-E & decision VOI vs.\ LLM proxy & 0.277 & $[0.119,\ 0.412]$ & 400 & 0.45/0.44/0.14 \\
    N4-F3 & typed transport vs.\ naive carry-forward & 17.242 & $[15.577,\ 18.836]$ & 200 & 0.69/0.30/0.00 \\
    N4-F3 & typed transport vs.\ re-derivation (mixed) & 2.264 & $[1.717,\ 2.825]$ & 200 & 0.32/0.69/0.00 \\
    N4-F3 & typed transport vs.\ re-derivation (no remint needed)$\ddagger$ & 0.000 & $[0.000,\ 0.000]$ & 200 & 0.00/1.00/0.00 \\
    \bottomrule
  \end{tabularx}
  \par\smallskip
  \textit{Metric note:} N4-A, N4-E and N4-F3 differences are in episode utility; N4-B in mean round utility; N4-C in scalarized regret reduction. N4-D is a frozen exact census (200 laundering / 200 honest chains), not a sampled paired comparison, and therefore has no interval; its counts are reported in the text.
\end{table}


\subsection{Negative first-refusal cases}
N1-C finds a typed/scoped failure-state benefit over a scoping ablation, but an ideal VOI allocator given the same typed facts matches the candidate policy exactly at solve rate 0.9866. The allocation-policy residual is therefore closed. N2-F5B similarly gives a model-selection donor first refusal: candidate and donor both score 0.9948 on the original world, leaving only the narrower misspecified-world edge (0.9844 versus 0.9531) as a surviving bounded result.